\tituloI{MARCO TEÓRICO}

\tituloII{Antecedentes de la investigación}
\parrafoII{
A nivel internacional, Smith y Brown \cite{smith2023global} destacan que los sistemas de gestión de inventarios modernos deben priorizar la validación de transacciones en el punto de origen para evitar discrepancias financieras. En el ámbito nacional, la SUNAT \cite{sunat2023informe} reporta que, aunque la masificación de los comprobantes electrónicos ha avanzado, persiste una brecha técnica en las microempresas que utilizan software no integrado, lo que García y Quispe \cite{garcia2023gestion} identifican como el principal factor de incumplimiento tributario involuntario en regiones como Cusco.
}

\tituloII{Bases teóricas}

\tituloIII{Transformación Digital en el Retail}
\parrafoIII{
La digitalización en el sector minorista implica la sustitución de procesos manuales por flujos de trabajo automatizados que aseguren la disponibilidad de información para la toma de decisiones estratégicas \cite{vargas2022transformacion}.
}

\tituloIII{Integridad de Datos y Gestión de Inventarios}
\parrafoIII{
La integridad referencial y las restricciones de dominio en bases de datos son fundamentales para prevenir el fenómeno del "stock negativo". Un sistema robusto debe garantizar que cada salida de almacén esté respaldada por una entrada previa verificada \cite{smith2023global}.
}

\tituloIII{Facturación Electrónica en el Perú}
\parrafoIII{
El modelo de facturación directa implica la generación de archivos XML firmados digitalmente bajo el estándar UBL (Universal Business Language) y su envío mediante servicios web SOAP a los servidores de la administración tributaria \cite{sunat2023informe}.
}

\tituloIII{Arquitectura de Sistemas Locales}
\parrafoIII{
En entornos de hardware limitado, la arquitectura de sistemas locales debe optimizar el uso de recursos de CPU y memoria, priorizando bases de datos ligeras pero potentes como MariaDB para mantener la operatividad "Bare Metal" sin depender de capas de virtualización costosas \cite{garcia2023gestion}.
}z